\documentclass[12pt]{aunarAPA}  

% cargar paquetes según la necesidad
\usepackage[spanish]{babel}
\usepackage[hidelinks]{hyperref}
\usepackage[utf8]{inputenc}
% \usepackage{showframe}
\usepackage{enumerate}

\usepackage{amssymb}
\usepackage{amsmath}
\usepackage{amsfonts}
\usepackage{mathrsfs}

\usepackage{graphicx}
\usepackage{times}
\usepackage{caption}
% \usepackage[caption=false,font=footnotesize]{subfig}
\usepackage{enumitem}
\usepackage{epstopdf}
\usepackage{amsthm}
\usepackage{bm}
\usepackage{cleveref}

% \usepackage{xcolor}
\usepackage{pgfplots}
\pgfplotsset{compat=newest}
\pgfplotsset{plot coordinates/math parser=false}
\usepackage{pgfgantt}
% \usepackage{rotating}
\usepackage{array}
\usepackage{multirow}
\usepackage{xcolor,colortbl}
% \usepackage[table]{xcolor}
\usepackage{pdfpages}
\usepackage{siunitx}
% \usepackage{biblatex}
%% paquete adicional para configuración en formato APA
\usepackage{aunarAPA}


%% Valores necesarios para crear la portada y subportada

\title{Titulo trabajo de grado} 

\author{Nombres y apellidos del autor} % Nombres y Apellidos del autor
% \author{Autor 1 \par Autor 2 \par Autor 3}% en el caso de que haya mas de un autor
\institucion{Corporación Universitaria Autónoma de Nariño}  %Universidad
\facultad{Ingeniería}
\programa{Ingeniería Electrónica}
\ciudad{San Juan de Pasto}
\fecha{15 de Enero}         % fecha de sustentación
\yeardef{2020}
\grado{ingeniero electrónico}
\asesortec{Nombre asesor técnico}
% \asesortec{Nombre asesor técnico 1 \\Nombre asesor técnico 2}
\asesormet{Nombre asesor metodológico}


\begin{document}

% Determina el numero de niveles para la tabla de contenido.
\setcounter{secnumdepth}{4}
\setcounter{tocdepth}{3}

\portada        % Crea la portada de la pagina
\begin{romanpages}          % Empieza numeración romana
    \subportada     % Crea la subportada del trabajo
    \aceptacion
    \include{Preambulos/dedicatoria}       % Incluye la dedicatoria
    \include{Preambulos/agradecimientos}   % Incluye los agradecimientos
    % \include{Capitulos/RAE}             % RAE
    \tableofcontents            % Tabla de contenido
    \listoffigures              % Lista de figuras
    \listoftables
\end{romanpages}            % Termina numeración romana
\setlength{\parindent}{1.27cm} % Sangría al inicio de cada parrafo

% Ahora incluya todo los capítulos y partes del trabajo
\include{Capitulos/Resumen}
\include{Capitulos/Introduccion}
\include{Capitulos/Planteamiento}
\include{Capitulos/Objetivos}

\include{Capitulos/MarcoReferencial}
\include{Capitulos/Metodologia}
\include{Capitulos/Administracion}
\include{Capitulos/Resultados}
\include{Capitulos/Conclusiones}

%Incluya apéndices si necesario

\appendix
\include{Anexos/anexoA}
% \end{appendices}

% Esta linea añade la bibliografía a la tabla de contenido
% \addcontentsline{toc}{chapter}{Lista de referencias}
\bibliography{Biblio}        %usar un archivo bibtex Biblio.bib
\bibliographystyle{apacite}  %Estilo APA


\end{document}